Today's denoising diffusion models do not ``denoise'' in the classical sense, \ie, they do not directly predict clean images. 
Rather, the neural networks predict noise or a noised quantity.
In this paper, we suggest that predicting clean data and predicting noised quantities are fundamentally different.
According to the manifold assumption, natural data should lie on a low-dimensional manifold, whereas noised quantities do not. With this assumption, we advocate for models that directly predict clean data, which allows apparently under-capacity networks to operate effectively in very high-dimensional spaces. We show that simple, large-patch Transformers on pixels can be strong generative models: using no tokenizer, no pre-training, and no extra loss. Our approach is conceptually nothing more than ``\mbox{\textbf{Just image Transformers}}'', or \textbf{JiT}, as we call it. We report competitive results using JiT with large patch sizes of 16 and 32 on ImageNet at resolutions of 256 and 512, where predicting high-dimensional noised quantities can fail catastrophically. 
With our networks mapping back to the basics of the manifold,
our research goes back to basics and pursues a self-contained paradigm for Transformer-based diffusion on raw natural data.
